\section{Limitations}
\label{sec:limitations}

\textbf{The CNN was not trained.} The analytics node runs a 1D-CNN with random weights, sized so one inference takes about 40\,ms on the board. It is realistic background load, not a fault classifier, so the setpoints sent to the MCU come from the FFT rule, not from the CNN. This does not affect the timing study, but we make no claim about fault-detection accuracy.

\textbf{The MCU loop does not use a Zephyr timer.} Our plan specified a periodic Zephyr \texttt{k\_timer}. All MCU results use the Arduino sketch's \texttt{micros()} schedule with a 6\,ms sleep and a final spin. A \texttt{k\_timer} version was never run on the board. The loop also has a latent flaw found in code review: if a cycle ever started more than about 3.5\,ms late, its 6\,ms sleep would overshoot the next release. This never happened (C's worst case was about 10\us), and the measured code was left unchanged.

\textbf{Several explanations are untested.} (i) Earlier single runs showed Linux slower at idle than under load (median 224\us{} for A and 171\us{} for B at idle, 44\us{} for B under load). We attribute this to waking from deep CPU idle states. It is consistent with B-iso, which disables them, having a 5$\times$ smaller idle worst case than B (96 vs 471\us), but B-iso also changes pinning and interrupt placement, so the effect is not isolated. (ii) In an early test, locking memory with \texttt{mlockall} under memory pressure gave \emph{higher} lateness than not locking it (mean 413 vs 208\us); this was not explained, and we do not cite it as evidence for or against \texttt{mlockall}. (iii) The kernel work on B-iso's core is consistent with its higher lateness but was not traced per cycle; a kernel trace would settle it.

\textbf{No \texttt{PREEMPT\_RT} kernel.} B and B-iso ran on the stock \texttt{PREEMPT} kernel. \texttt{PREEMPT\_RT} moves most interrupt and softirq work into schedulable threads~\cite{reghenzani2019rtlinux} and would likely narrow the gap between B and C; our Linux results are therefore not the best real-time Linux can do.

\textbf{Scope of the evidence.} One board set; an emulated drive instead of a motor; two repeats per realistic condition, the two \texttt{rboth} repeats in one session. MCU runs are 18$\times$ shorter than Linux runs, which is why we compare per block and why C's zero misses bound its miss ratio only at 0.015\,\%. The CNN runtime lets idle worker threads spin by default~\cite{onnxruntime_threads}; we kept the default, as a typical deployment would, so the CNN tiers occupy more CPU than inference needs.
